\section{Agents, Baselines and Protocol}
\label{sec:method}

\subsection{Learning Agents and Tuning}
Three deep RL algorithms are trained. PPO \cite{schulman2017} and A2C \cite{mnih2016} learn a policy and a value function. DQN \cite{mnih2015} learns action values. All three come from Stable-Baselines3 \cite{raffin2021} and use the Gymnasium interface \cite{towers2024}. Agents are tested with their greedy action.

Every algorithm gets the same tuning budget, so that no reward or algorithm is helped by better settings. For each algorithm, six configurations are tried: the setting of the conference version and five drawn at random from a fixed grid (Table~\ref{tab:hyper}). For DQN the grid includes multiplying all rewards by 0.1, which does not change the optimal policy. Each configuration is trained with both the weighted and the flat reward, with two seeds (100 and 101) that are not used later, for 50,000 steps. The configuration with the highest exact expected return under its own reward is chosen. The SLA cost is never used for this choice, so the tuning does not favor the operator's objective. The corrected reward uses the setting chosen for the weighted reward.

\begin{table}[t]
\caption{Tuning grid and chosen settings. The same setting was chosen for both rewards for every algorithm. Unlisted settings are Stable-Baselines3 2.9.0 defaults.}
\label{tab:hyper}
\centering
\footnotesize
\setlength{\tabcolsep}{2.5pt}
\begin{tabular}{@{}l>{\raggedright\arraybackslash}p{3.6cm}l@{}}
\toprule
Algorithm & Grid & Chosen \\
\midrule
PPO & learning rate $\{1,3,10\}\times10^{-4}$, steps per update $\{512, 1024, 2048\}$, batch $\{64,128\}$, entropy $\{0, 0.01\}$, $\gamma$ $\{0.99, 0.999\}$, layers $\{2\times64, 2\times128\}$ & PPOCHOSEN \\
\addlinespace
DQN & learning rate $\{1,3,10\}\times10^{-4}$, target update $\{500, 2000\}$, exploration share $\{0.1,0.3\}$, buffer $\{10^5,10^6\}$, $\gamma$ $\{0.99,0.999\}$, reward scale $\{1, 0.1\}$, layers $\{2\times64, 2\times128\}$ & DQNCHOSEN \\
\addlinespace
A2C & learning rate $\{3,7,10\}\times10^{-4}$, steps per update $\{5, 32, 128\}$, $\gamma$ $\{0.99,0.999\}$, entropy $\{0,0.01\}$, normalized advantage $\{\text{no},\text{yes}\}$, layers $\{2\times64, 2\times128\}$ & A2CCHOSEN \\
\bottomrule
\end{tabular}
\end{table}

With the chosen settings, each agent is trained with 10 seeds, 0 to 9, for 100,000 steps in the nominal system. The exact model scores every agent every 5,000 steps, which gives learning curves without test noise. In the larger systems PPO is trained for 100,000 steps in Hetero3 and Hetero3-3C and 200,000 steps in Large4-3C, again with 10 seeds and all three rewards. Because the rules are tuned on the exact model of each system, PPO is tuned for each system too. The three PPO settings with the highest mean return in the nominal tuning are trained with both rewards, the two tuning seeds and the full budget of the system, and the one with the highest exact expected return is chosen. All 10-seed runs are also repeated with the nominal PPO setting.

\subsection{Rules}
Four rules serve as baselines, and two of them are tuned on the exact model of each system. Join-shortest-queue (JSQ) picks the edge node with the fewest waiting tasks and never uses the cloud. The SLA rule knows Eq.~\eqref{eq:latency} and every limit. It picks the edge node with the lowest predicted latency and uses it if that latency meets the limit of the task, and the cloud otherwise. The reservation rule adds one latency threshold for each class except the most critical one. A task of such a class goes to the edge only if the predicted latency of the best node is at most its threshold, which keeps capacity free for critical tasks. This is the form of the optimal admission policy for one pool of equal servers \cite{miller1969}. The thresholds are chosen by exhaustive search for two classes, and by coordinate search started from the SLA rule for three classes, minimizing the exact SLA cost. The reservation rule with node sets, called the node-set rule below, also lets each class use only a chosen subset of the edge nodes, for example only the slow nodes for analytics tasks. Subsets and thresholds are chosen by coordinate search on the exact model, started from the tuned reservation rule. All rules use the true latency model and the exact model of the system, which no learning agent sees. In the nominal system the node sets bring no gain, because the two nodes are equal.

\subsection{Published Reward Designs}
The flat reward is the design of Nieto et al. \cite{nieto2024}, in which the class is part of the state but every failed task has the same penalty. The design of Yamsani and Chenna Reddy \cite{yamsani2026} is rebuilt in three parts, following their equations. First, the penalty of class $c$ starts at $e_c$ and is updated after every task as $\lambda_c \leftarrow \max(0, \lambda_c + \eta\,(V_c - V^\ast))$, where $V_c$ is the miss rate of class $c$ over its last 200 tasks and $V^\ast = 0.05$ is their 95\% compliance target. Second, the state is extended with risk features that carry the same information as their risk score, which combines the expected latency against the deadline, the queue length and the recent miss history. The features are the predicted latency of each edge node divided by the limit of the task and the recent miss rate of each class. Third, following their pruning of clearly unsafe actions, edge nodes whose predicted latency misses the limit of the task are removed from the action set by invalid-action masking \cite{huang2022}. Their paper does not give the step size $\eta$, so two values are tested, 0.001 and 0.01. The agent is PPO with the tuned PPO setting, so that the comparison isolates the design and not the learner. A second version keeps only the adaptive penalties, without risk features and pruning. The full design uses features that are not part of the state of Section~\ref{sec:steps}, so it is scored by simulating 500 held-out episodes. All other agents are scored exactly.

\subsection{Changes of the System}
Five changes of the nominal system test what happens when the system differs from the one that rules and agents were built for. Tasks get slower, with 20\,ms instead of 15\,ms per waiting task. Service gets slower, with $d = 0.55$. Background traffic rises to 0.3 tasks per step. Background traffic becomes random with the same mean. The cloud meets the analytics limit with probability 0.5. For each change, four kinds of policies are scored exactly. The rules keep their nominal latency model and threshold. The trained PPO agents keep their nominal policy. The rules are tuned again with the true model of the changed system. The PPO agents are trained for another 25,000 steps in the changed system, starting from their nominal weights, which uses only interaction and rewards. The optimum of the changed system gives the reference.

\subsection{Statistics}
For learning agents the paper reports the mean over seeds with a 95\% bootstrap confidence interval from 10,000 resamples \cite{efron1993}. Rules and DP policies are deterministic and evaluated exactly, so they have no interval. Two groups of agents are compared with the two-sided Mann-Whitney U test \cite{mann1947}, which is exact when the samples have no tied values. The effect size is the Vargha-Delaney $A_{12}$ \cite{vargha2000}, the probability that a seed of the first group has a higher value than a seed of the second. Tests within one table are corrected with the Holm method \cite{holm1979}. Agents are compared with a deterministic rule by a two-sided sign test over seeds.
